% TOP_PROBLEM: {"id": 3292, "title": "Ryser's conjecture", "category_id": 3, "category": {"id": 3, "name": "graph_theory", "display_name": "Graph Theory", "description": "Problems involving graphs, networks, and their properties.", "slug": "graph-theory", "order_index": 3, "created_at": "2026-07-31T15:26:25.671Z"}, "status": "open", "problem_number": "OPG-165", "set_id": 12}
% FIRST_POSTED: 2026-10-01
% ENTRY_KIND: statement_only
% UnsolvedMath ID 3292; source catalogue code OPG-165
% !TeX program = lualatex
% Definition and sources only; this file is not an LLM solution attempt.
\documentclass[11pt,a4paper]{article}
\usepackage[margin=25mm]{geometry}
\usepackage{fontspec}
\setmainfont{lmroman10-regular.otf}[BoldFont=lmroman10-bold.otf,
 ItalicFont=lmroman10-italic.otf,BoldItalicFont=lmroman10-bolditalic.otf]
\usepackage{amsmath,amssymb,mathtools}
\usepackage{unicode-math}
\setmathfont{latinmodern-math.otf}
\AtBeginDocument{\let\setminus\smallsetminus}
\usepackage{xurl,hyperref}
\hypersetup{colorlinks=true,urlcolor=blue}
\setcounter{secnumdepth}{0}
\setlength{\parindent}{0pt}
\setlength{\parskip}{5pt}
\setlength{\emergencystretch}{3em}
\begin{document}
\section*{Ryser's conjecture}\label{3292}
{\small UnsolvedMath ID 3292 · OPG-165 · Graph Theory · OpenGarden}
\subsection{Definitions and mathematical statement}
Fix an integer $r\ge2$. A finite $r$-uniform hypergraph $H=(V,E)$ has
edges which are $r$-element subsets of $V$. It is $r$-partite if there
is a partition $V=V_1\sqcup\cdots\sqcup V_r$ such that every edge
contains exactly one vertex of each part. The parts may have different
sizes. Repeated edges are unnecessary and are not included.

A matching is a family of edges with pairwise disjoint vertex sets;
let $\nu(H)$ be its maximum possible size. A vertex cover, or
transversal, is a set $S\subseteq V$ with $S\cap e\ne\varnothing$ for
every $e\in E$; let $\tau(H)$ be the minimum size of such a set. If
there are no edges both parameters are zero.

Ryser's conjecture asks whether every finite $r$-partite $r$-uniform
hypergraph satisfies
\[
                             \tau(H)\le(r-1)\nu(H).
\]
There is no regularity or bound on degrees. Covers may use vertices
from several parts, and are not required to lie in one part.

The inequalities $\nu(H)\le\tau(H)\le r\nu(H)$ give the basic scale.
Each disjoint edge needs a separate cover vertex, while the union of
a maximum matching meets every edge and contains $r\nu(H)$ vertices.
The conjecture saves a whole matching-sized portion by exploiting the
partite structure. For $r=2$ its conclusion is equality of the matching
and cover numbers of a bipartite graph. The restriction $r\ge2$
avoids the false coefficient-zero assertion for nonempty 1-uniform
hypergraphs.
\subsection{Short English statement}
In an r-partite r-uniform hypergraph, is the minimum vertex-cover size at most r−1 times the maximum matching size?
\subsection{Sources}
\begin{itemize}
\item[S1] ULAM.AI and the UnsolvedMath Contributors, supplied problems.json, record 3292 (OPG-165), OpenGarden collection. Catalogue identity and source transcription; CC BY 4.0.
\url{https://huggingface.co/datasets/ulamai/UnsolvedMath}.
\item[S2] Open Problem Garden, Ryser's conjecture. Original problem and discussion; the mathematical formulation below is expanded from the supplied source record.
\url{http://www.openproblemgarden.org/op/rysers_conjecture}.
\end{itemize}
\end{document}
